\section{The Mathematician, the Carpenter, and the Physicist}

\begin{quotex}
Mathematics, rightly viewed, possesses not only truth, but supreme beauty cold and austere, like that of sculpture, without appeal to any part of our weaker nature, without the gorgeous trappings of painting or music, yet sublimely pure, and capable of a stern perfection such as only the greatest art can show. The true spirit of delight, the exaltation, the sense of being more than Man, which is the touchstone of the highest excellence, is to be found in mathematics as surely as in poetry.
\flright{\textsc{Bertrand Russell}, \emph{Study of Mathematics}}
\end{quotex}


A mathematician, a carpenter, and a physicist walk into a bar. The carpenter is touching and admiring the bar's woodwork. The mathematician is siting and thinking, while the physicist was busily doodling on his napkin. After a while, they struck up a conversation.

The carpenter asked the mathematician, ``I have a client who wants a square herb garden, of 4 square feet. How long should I make the sides?''

The mathematician responded, ``There are two square roots of 4: +2 and -2. That is how long each side should be.''

The carpenter seemed confused and thought to himself, ``I don't know how to cut a piece of wood of length -2 feet, so I'll just cut 4 2-foot lengths of wood and use them to shape the garden.'' He thanked the mathematician and left.

The physicist was eagerly eavesdropping and doodled all the more. ``Look,'' he said to the mathematician, barely able to conceal his excitement, ``if I plug -2 into my equations, I can create another parallel world to ours.'' Then he rushed out, in search of a large blackboard.

The mathematician kept thinking, because any practical applications of his maths did not matter to him.

\paragraph{The Bartender}


The bartender, who was actually a metaphysician who could not get a job in that field, seemed amused at what transpired. As he wiped the counter, he struck up a conversation with the mathematician. He said, ``I noticed that you gave the physicist a universal principle which he then applied to his physical theory.''

``Yes, and after he does his experiments and publishes his theory, he will claim to have found a law of the universe,'' replied the mathematician, rather nonchalantly.

The bartender chuckled. ``Then the academic philosophers will chortle about the unreasonable effectiveness of mathematics in the natural sciences\footnote{\url{https://www.gornahoor.net/library/UnreasonableEffectivenessOfMath.pdf}}. Mathematics is the dividing line between metaphysics and physics. As such, it is the principle of individuation through which the multifarious empirical phenomena emerge.''

The other fellows left without paying, so the mathematician paid the tab, leaving a nice gratuity to the bartender.


\flrightit{Posted on 2021-09-19 by Cologero}

